\begin{enumerate}[leftmargin=*,itemsep=0.55em]
\item La base organizada cumple con el número de registros solicitado: diez compuestos originales, cien reacciones y cien obtenciones. La auditoría confirma la preservación de los campos descriptivos cotejados y la concordancia de 1.060 conversiones numéricas simples, sin garantizar por ello la exactitud fisicoquímica de todas las fuentes.
\item Las categorías más observadas en los registros son oxidación (17,5\%), hidrogenación (13,0\%) y reducción (11,0\%). Su distribución depende de los criterios de selección y del rol reacción/obtención; no debe extrapolarse a la abundancia de rutas existentes en CAS-SciFinder.
\item El medio neutro es la clase más frecuente (38,5\%), aunque el 27,5\% de registros se clasifica como «otro/no determinado». Este porcentaje limita cualquier interpretación general sobre el papel de la acidez o basicidad.
\item La mediana de la temperatura central es 100~$^\circ$C para los 170 registros con valores numéricos, pero los registros multietapa y las condiciones no homogéneas impiden interpretar dicho valor como una temperatura característica de la reactividad de alcanos o alquenos.
\item La comparación butano--isobutano muestra una diferencia de 11,2~$^\circ$C en el punto de ebullición tabulado a masa molecular idéntica. El ejemplo es compatible con la importancia de la ramificación en las interacciones intermoleculares, aunque no permite inferir causalidad estadística.
\item La separación entre compuesto base, producto principal y reactante evaluado fue esencial para evitar mezclar propiedades de especies con roles diferentes. La cobertura de los datos se incorporó al análisis para interpretar las comparaciones con el denominador correspondiente.
\end{enumerate}


